\documentclass[12pt,a4paper]{amsart}

\usepackage{amsmath,amssymb,amsthm}
\usepackage{mathtools}
\usepackage[colorlinks=true,linkcolor=blue!60!black,citecolor=green!50!black,
            urlcolor=blue!70!black]{hyperref}
\usepackage[shortlabels]{enumitem}
\usepackage{booktabs}

%% ---- Theorem environments ----------------------------------------
\newtheorem{theorem}{Theorem}[section]
\newtheorem{lemma}[theorem]{Lemma}
\newtheorem{proposition}[theorem]{Proposition}
\newtheorem{corollary}[theorem]{Corollary}
\theoremstyle{definition}
\newtheorem{definition}[theorem]{Definition}
\newtheorem{remark}[theorem]{Remark}
\newtheorem{assumption}[theorem]{Assumption}
\newtheorem{conjecture}[theorem]{Conjecture}
\newtheorem{problem}{Problem}[section]

%% ---- Notation (aligned with Paper I) ----------------------------
\newcommand{\PW}{\mathrm{PW}_\omega}
\newcommand{\Kop}{\mathcal{K}_c}
\newcommand{\GVM}{G^{(N)}_{\mathbf{p},c}}
\newcommand{\QM}{Q^{(N)}_{\mathbf{p},c}}
\newcommand{\psin}{\psi_n^{(c)}}
\newcommand{\EN}{E_{N,c}}
\newcommand{\norm}[1]{\|#1\|}
\newcommand{\ip}[2]{\langle #1,\, #2\rangle}
\newcommand{\Rbulk}{\mathcal{R}_{\mathrm{bulk}}}
\newcommand{\Rturning}{\mathcal{R}_{\mathrm{turn}}}
\newcommand{\PN}{P_{N,c}}
\newcommand{\CXRY}{C_{\mathrm{XRY}}}

\title[Conditional Frame Stability via PSWF Spectral Projection]
      {Conditional Frame Stability of Prolate Sampling Operators\\
       via PSWF Spectral Projection and XRY Quadrature:\\
       A Conditional Implication Framework for Quadrature Compatibility}

\thanks{This paper is a companion to \cite{PaperI}, where frame coercivity
and a conditional banded decay estimate were established under an explicit
quadrature compatibility assumption (Assumption~3.10 of \cite{PaperI}).
The present paper does not prove that assumption outright. Instead, it
organises the problem into two explicit conjectural components: a PSWF
product spectral tail conjecture and an XRY spectral stability conjecture.
Their combination yields a conditional verification mechanism for the
compatibility assumption used in \cite{PaperI}.}

\author{Anonymous Author}
\address{}
\email{}
\date{\today}

\subjclass[2020]{42C05, 42B10, 47B32, 41A60, 65D32}

\keywords{prolate spheroidal wave functions, PSWF-Gauss quadrature,
finite frames, frame stability, Gram operator, spectral projection,
quadrature compatibility, Xiao--Rokhlin--Yarvin, exponential convergence}

\begin{document}
\maketitle

\begin{abstract}
In the companion paper \cite{PaperI}, frame coercivity and a conditional
banded decay estimate for the quadrature defect matrix $\EN$ were
established for the prolate sampling operator on $V_{N,c}$, conditional
on an explicit quadrature compatibility assumption (Assumption~3.10 of
\cite{PaperI}).

The present paper does not verify that assumption unconditionally.
Instead, for PSWF-Gauss quadrature rules of Xiao--Rokhlin--Yarvin
\cite{Xiao2001}, it organises the compatibility problem into two explicit
conjectural inputs. First, one needs a PSWF product spectral tail statement
for the analytic product family
$f_{mn}=\psi_m^{(c)}\overline{\psi_n^{(c)}}$ in the relevant index regime.
Second, one needs an XRY spectral stability statement on the same family.
Assuming these two conjectural inputs, one obtains the conditional defect
estimate
\[
  |E_{mn}| \le C e^{-\alpha N} \qquad (0\le m,n\le N-1).
\]
Therefore Assumption~3.10 of \cite{PaperI} is verified conditionally, and
the frame stability theorem of \cite{PaperI} applies to PSWF-Gauss sampling.
\end{abstract}

\tableofcontents

%% ==================================================================
\section{Introduction}
%% ==================================================================

\subsection*{Context}

The companion paper \cite{PaperI} established the following structure
for the quadrature defect matrix $\EN = \GVM - G^{\mathrm{ref}}_{N,c}$:
\begin{enumerate}
\item An exact identity $E_{mn} = R_{mn}^{\mathrm{quad}}$
  (pure quadrature remainder; PSWF orthogonality term vanishes exactly).
\item A conditional banded decay estimate: under Assumption~3.10 of
  \cite{PaperI},
  \[
    |E_{mn}| \le C_p(1+|m-n|)^{-p} + Cc^{-1/3}
    \qquad (0 \le m,n \le N-1).
  \]
\item A conditional frame stability theorem with explicit bounds.
\end{enumerate}
The open problem left by \cite{PaperI} was to identify a concrete
quadrature class satisfying Assumption~3.10.

\subsection*{Contribution of this paper}

This paper does not solve the full construction problem unconditionally.
Its contribution is an implication framework: for XRY PSWF-Gauss quadrature,
Assumption~3.10 of \cite{PaperI} follows from two separate conjectural
inputs, namely a PSWF product spectral tail conjecture and an XRY spectral
stability conjecture on the relevant product family
$f_{mn}=\psi_m^{(c)}\overline{\psi_n^{(c)}}$.

\subsection*{Strategy of this paper}

The technical strategy consists of an implication chain with two explicit
analytical inputs:
\begin{enumerate}[(i)]
\item \emph{PSWF product spectral tail conjecture.} The product integrands
  $f_{mn}=\psi_m^{(c)}\overline{\psi_n^{(c)}}$ are expected to admit
  sufficiently uniform PSWF spectral tail decay in the regime
  $0\le m,n\le N\sim c$.
\item \emph{XRY spectral stability conjecture.} The XRY quadrature
  functional $Q_N$ is expected to satisfy a stability bound on the same
  product family, controlled by the PSWF spectral truncation error
  associated with $\PN$.
\end{enumerate}
Within the present framework, these two inputs are used only through their
combination in the deduction of the defect estimate. The technical content
of the paper is therefore the exact implication chain from these two
conjectural inputs to the compatibility condition of \cite{PaperI}, together
with the perturbative passage from that condition to frame stability.

\subsection*{Main result}

\begin{theorem}[Conditional implication theorem]\label{thm:main_informal}
Assume Conjectures~\ref{conj:pswf_product_tail} and~\ref{conj:xry_stability},
together with the auxiliary technical hypotheses stated in
Assumptions~\ref{ass:pswf_sup}, \ref{ass:node_sep}, and~\ref{ass:weights}.
Then for PSWF-Gauss quadrature rules \cite{Xiao2001}, the quadrature defect
satisfies $|E_{mn}|\le Ce^{-\alpha N}$ for all $0\le m,n\le N-1$.
Consequently, Assumption~3.10 of \cite{PaperI} holds conditionally on the
same inputs, and the frame stability theorem of \cite{PaperI} applies to
PSWF-Gauss sampling.
\end{theorem}

The formal version is Theorem~\ref{thm:compatibility_verified} together
with Theorem~\ref{thm:frame_stability}.

\subsection*{Positioning}

This paper should be read as an implication framework (in the sense of a
reduction between explicit analytic hypotheses) rather than a Theorem-paper
in the traditional sense. Its purpose is to isolate the two core unresolved
analytical components that stand between the abstract framework of
\cite{PaperI} and a concrete verification for XRY PSWF-Gauss quadrature.
The PSWF product-tail step and the XRY stability step are both left as
explicit conjectural inputs. In particular, the present work is not a
classical analysis paper deriving the conclusion from weaker previously
established estimates; rather, it makes explicit where the unresolved
analytic burden lies and how the downstream frame-theoretic conclusion would
follow once those inputs are available.

\subsection*{Organisation}

Section~\ref{sec:pswf_gauss} sets up PSWF-Gauss quadrature and the
spectral projector. Section~\ref{sec:spectral_decay} states the PSWF-side
conjectural input and explains the scope of the product-tail step.
Section~\ref{sec:xry_accuracy} states the XRY spectral stability conjecture.
Section~\ref{sec:verification} derives the conditional verification of
Assumption~3.10. Section~\ref{sec:stability} states the conditional frame
stability theorem. Section~\ref{sec:open} records open problems.

%% ==================================================================
\section{PSWF-Gauss Quadrature and Spectral Projector}\label{sec:pswf_gauss}
%% ==================================================================

\begin{definition}[PSWF spectral projector]\label{def:projector}
Let $\{\psi_k^{(c)}\}_{k\ge0}$ be an orthonormal PSWF basis in
$L^2([-T,T])$. For $N\ge1$, define
\[
  \PN f := \sum_{k=0}^{N-1}\ip{f}{\psi_k^{(c)}}\,\psi_k^{(c)}.
\]
\end{definition}

\begin{definition}[PSWF-Gauss quadrature operator]\label{def:xry}
Let $(\mathbf{p},\mathbf{w})=((p_1,\dots,p_N),(w_1,\dots,w_N))$ denote
the PSWF-Gauss nodes and weights of Xiao--Rokhlin--Yarvin \cite{Xiao2001}.
The quadrature functional is $Q_N(f):=\sum_{j=1}^N w_j f(p_j)$.
\end{definition}

\begin{remark}[No algebraic closure claim]\label{rem:no_algebra}
We do not assume any closure of the PSWF system under pointwise
multiplication. All arguments below use only analytic regularity and
spectral approximation properties.
\end{remark}

\begin{assumption}[Node separation]\label{ass:node_sep}
The PSWF-Gauss nodes satisfy
$\min_{j\ne k}|p_j-p_k|\ge c_0 T/N$
for some constant $c_0>0$ in the regime $N\sim c$.
\end{assumption}

\begin{assumption}[Weight positivity and boundedness]\label{ass:weights}
The weights satisfy $w_j>0$ and $C_1 T/N\le w_j\le C_2 T/N$
uniformly in $j$ in the regime $N\sim c$.
\end{assumption}

\begin{remark}[Node-weight assumptions and evidence]\label{rem:nodes_weights_evidence}
Assumptions~\ref{ass:node_sep} and~\ref{ass:weights} are included here as
explicit hypotheses rather than derived facts. They are, however,
consistent with the numerical behaviour reported for PSWF-Gauss rules in
\cite{Xiao2001,OsipovRokhlinXiao}, where the nodes appear well separated and
the weights remain positive and of the expected order $T/N$ in the regime
$N\sim c$. We do not claim these properties as proved here.
\end{remark}

%% ==================================================================
\section{PSWF Product Spectral Tail Conjecture}\label{sec:spectral_decay}
%% ==================================================================

\begin{assumption}[Uniform PSWF sup-norm control]\label{ass:pswf_sup}
There exists a constant $C_{\infty}>0$ such that
\[
  \norm{\psi_n^{(c)}}_{L^\infty([-T,T])}\le C_{\infty}
\]
uniformly for $0\le n\le N\sim c$ in the bulk regime; compare the standard
PSWF bounds discussed in \cite{Slepian1978,OsipovRokhlinXiao}.
\end{assumption}

\begin{lemma}[Uniform product norm bound]\label{lem:product_norm}
Assume Assumption~\ref{ass:pswf_sup} and the regime $0\le m,n\le N\sim c$.
Then there exists a constant $C>0$ such that
\[
  \norm{f_{mn}}_{L^2([-T,T])}
  = \norm{\psi_m^{(c)}\overline{\psi_n^{(c)}}}_{L^2([-T,T])}
  \le C
\]
uniformly in $m,n$.
\end{lemma}
\begin{proof}
By Assumption~\ref{ass:pswf_sup},
$\norm{\psi_m^{(c)}}_{L^\infty([-T,T])}\le C_{\infty}$ uniformly in the
bulk regime, while $\norm{\psi_n^{(c)}}_{L^2([-T,T])}=1$ by orthonormality.
Hence
\[
  \norm{f_{mn}}_{L^2}
  \le \norm{\psi_m^{(c)}}_{L^\infty}\norm{\psi_n^{(c)}}_{L^2}
  \le C_{\infty}.
\]
\end{proof}

\begin{remark}[Role of the product norm bound]\label{rem:generic_bound}
The preceding estimate uses only the generic $L^\infty\cdot L^2$ bound and
carries no PSWF-specific spectral information. Its role is solely to ensure
that the $L^2$ factor in Conjecture~\ref{conj:xry_stability} remains
uniformly bounded; the nontrivial spectral content is entirely contained in
Conjecture~\ref{conj:pswf_product_tail}.
\end{remark}

\begin{remark}[Why the PSWF product-tail issue is difficult]
The family $f_{mn}=\psi_m^{(c)}\overline{\psi_n^{(c)}}$ is not known to enjoy
any finite-dimensional closure property in the PSWF basis. Thus a uniform
spectral-tail statement for $0\le m,n\le N\sim c$ is not a routine corollary
of standard PSWF approximation results, but a separate open analytical
problem.
\end{remark}

\begin{conjecture}[PSWF product spectral tail conjecture]\label{conj:pswf_product_tail}
For the family
\[
  \mathcal{F}_N:=\{\psi_m^{(c)}\overline{\psi_n^{(c)}}:0\le m,n\le N\},
\]
the following hierarchy is expected to hold in the regime $N\sim c$:
\begin{enumerate}[(i)]
\item \emph{Off-diagonal control.} For $|m-n|\gg 1$, the tails
  $\|(I-\PN)f_{mn}\|_{L^2([-T,T])}$ satisfy an exponentially decaying bound
  with constants uniform away from the diagonal.
\item \emph{Bulk control.} For $m,n\le \gamma N$ with fixed $\gamma<1$,
  the same bound holds with constants uniform in the bulk.
\item \emph{Uniform control up to the relevant index range.} There exist
  constants $C,\alpha>0$ and a sequence $\varepsilon_N\downarrow0$ such that
  for every $f_{mn}\in\mathcal{F}_N$,
  \[
    \|(I-\PN)f_{mn}\|_{L^2([-T,T])}
    \le C e^{-(\alpha-\varepsilon_N)N}
    \qquad(0\le m,n\le N\sim c).
  \]
\end{enumerate}
The third statement is the form actually used in the sequel. The preceding
staging is included to reflect a possible route toward verification and to
avoid suggesting that a fully uniform bound is the only reasonable entry
point. The strongest and most vulnerable part of this conjecture is the
uniformity near the edge regime $m,n\sim N$, where PSWF regularity is
expected to deteriorate and where any proof would have to control possible
collapse of effective analyticity widths.
\end{conjecture}

\begin{remark}[Frequency-space heuristic for the product-tail problem]
PSWFs with bandlimit parameter $c$ may be viewed heuristically as functions
whose Fourier content is essentially concentrated on a frequency window of
size comparable to $c$. Under this heuristic, the pointwise product
$\psi_m^{(c)}\overline{\psi_n^{(c)}}$ corresponds in frequency space to a
convolution, so one expects an effective bandwidth of order $2c$ rather
than $c$. Since the projector $\PN$ retains only the first $N\sim c$ PSWF
modes, the main issue is not whether spectral mass extends beyond that
range, but how rapidly the corresponding PSWF coefficients decay past the
cutoff. This is precisely where the edge regime $m,n\sim N$ becomes
critical.
\end{remark}

\begin{remark}[Status of intermediate structure]
To the best of our knowledge, there is currently no established
intermediate property of the PSWF product family that would serve as a
quantitative precursor to Conjecture~\ref{conj:pswf_product_tail}. In
particular, there is no known stable multiplicative closure scale in the
PSWF basis that would bridge the gap between generic analyticity and the
uniform tail estimate required here.
\end{remark}

\begin{remark}[Analytical bottleneck]\label{rem:bottleneck}
The overall implication chain shows that the main unresolved analytical
difficulty is concentrated in the PSWF product spectral tail, especially
near the edge regime $m,n\sim N$. Once such a tail estimate is available,
the remaining steps reduce to quadrature stability and perturbative frame
analysis.
\end{remark}

%% ==================================================================
\section{XRY Spectral Stability Conjecture}\label{sec:xry_accuracy}
%% ==================================================================

\begin{remark}[Why the XRY stability conjecture is structurally plausible]
The XRY quadrature rule is constructed independently through the
PSWF-Gauss node-weight scheme of Xiao--Rokhlin--Yarvin \cite{Xiao2001}, and
not by building the projector $\PN$ into the quadrature functional itself.
The role of $\PN$ in Conjecture~\ref{conj:xry_stability} is therefore
an a posteriori analytical comparison device rather than part of the
construction of $Q_N$. Moreover, the XRY formulas are exact on the
underlying PSWF trial space, so $Q_N$ integrates $\PN f$ exactly whenever
$\PN f\in V_{N,c}$. This suggests that $Q_N$ behaves as an approximate
spectral projector onto $V_{N,c}$ at the level of integration functionals,
with the defect governed by the contribution of $(I-\PN)f$. With that
distinction in place, it is reasonable to expect the quadrature error to be
governed primarily by the PSWF spectral tail of the test function, up to a
residual term that is exponentially small in the regime $N\sim c$. This is
also consistent with the interpolation formulas underlying the XRY
construction. This does not prove Conjecture~\ref{conj:xry_stability}, but
it explains why a bound of the form
\[
  |Q_N(f)-\int_{-T}^T f|\lesssim \|(I-\PN)f\| + e^{-\alpha N}\norm{f}_{L^2}
\]
represents a natural stability statement rather than an ad hoc
reformulation of the original compatibility problem.
\end{remark}

\begin{conjecture}[XRY stability conjecture on the PSWF product family]
\label{conj:xry_stability}
Let
\[
  \mathcal{F}_N:=\{\psi_m^{(c)}\overline{\psi_n^{(c)}}:0\le m,n\le N\}.
\]
There exist constants $\CXRY,C,\alpha>0$, depending on the XRY
construction parameters and on the regime $N\sim c$, such that for every
$f\in\mathcal{F}_N$,
\[
  \left|Q_N(f)-\int_{-T}^T f(x)\,dx\right|
  \le \CXRY\,\|(I-\PN)f\|_{L^2([-T,T])}
      + C\,e^{-\alpha N}\,\norm{f}_{L^2([-T,T])}.
\]
This conjecture formalises the XRY-side stability input needed in the
sequel. It is consistent with the structural formulas and numerical
behaviour of \cite{Xiao2001,OsipovRokhlinXiao}, but is not claimed here as a
known theorem.
\end{conjecture}

\begin{remark}[Why this is an implication framework]
Conjectures~\ref{conj:pswf_product_tail} and~\ref{conj:xry_stability}
isolate the two unresolved analytical inputs used in this paper.
Accordingly, the present work should be read as an implication framework
for the quadrature compatibility problem of \cite{PaperI}, not as an
unconditional verification of that compatibility condition.
\end{remark}

%% ==================================================================
\section{Conditional Verification of Assumption~3.10}\label{sec:verification}
%% ==================================================================

\begin{proposition}[Structured defect bound]\label{prop:structured_defect}
Assume Conjecture~\ref{conj:xry_stability}. Then for all $0\le m,n\le N-1$,
\[
  |E_{mn}|
  \le \left|Q_N(f_{mn})-\int_{-T}^T f_{mn}(x)\,dx\right|
  \le \CXRY\,\|(I-\PN)f_{mn}\|_{L^2([-T,T])}
      + C\,e^{-\alpha N}\,\norm{f_{mn}}_{L^2([-T,T])}.
\]
\end{proposition}
\begin{proof}
The first inequality is the exact identity
$E_{mn}=R_{mn}^{\mathrm{quad}}=Q_N(f_{mn})-\int_{-T}^T f_{mn}(x)\,dx$
from \cite{PaperI}. The second inequality is precisely the statement of
Conjecture~\ref{conj:xry_stability} applied to $f=f_{mn}$.
\end{proof}

\begin{proposition}[Conditional exponential quadrature defect bound]\label{prop:defect_bound}
Assume Conjectures~\ref{conj:pswf_product_tail}
and~\ref{conj:xry_stability}, together with Assumption~\ref{ass:pswf_sup}.
Then for all $0\le m,n\le N-1$,
\[
  |E_{mn}|=|R_{mn}^{\mathrm{quad}}|\le C\,e^{-\alpha N},
\]
where $C,\alpha>0$ depend on the conjectural constants and on the product
norm bound from Assumption~\ref{ass:pswf_sup}.
\end{proposition}
\begin{proof}
By Proposition~\ref{prop:structured_defect},
\[
  |E_{mn}|\le \CXRY\|(I-\PN)f_{mn}\|+Ce^{-\alpha N}\norm{f_{mn}}_{L^2}.
\]
By Conjecture~\ref{conj:pswf_product_tail}(iii),
$\|(I-\PN)f_{mn}\|\le C' e^{-(\alpha'-\varepsilon_N) N}$
uniformly over $0\le m,n\le N\sim c$, and by
Lemma~\ref{lem:product_norm}, $\norm{f_{mn}}_{L^2}\le C$ uniformly in the
same regime. All exponential contributions are therefore absorbed into a
single rate $\alpha$, possibly reduced, and all multiplicative constants are
absorbed into $C$. This yields the stated bound.
\end{proof}

\begin{remark}[No gradual interpolation in the present framework]
Within the present argument, there is no gradual weakening of the
conclusion obtainable from partial versions of the two conjectures. The
exponential defect bound follows only once both conjectural inputs are
available in the stated form.
\end{remark}

\begin{proposition}[Partial compatibility under bulk/off-diagonal control]
\label{prop:partial_compatibility}
Assume Conjecture~\ref{conj:xry_stability}, Assumption~\ref{ass:pswf_sup},
and the following partial PSWF tail information:
\begin{enumerate}[(a)]
\item for indices in the bulk regime $m,n\le \gamma N$ with fixed
  $\gamma<1$, one has a uniform algebraic estimate
  \[
    \|(I-\PN)f_{mn}\|_{L^2([-T,T])}
    \le C_p(1+|m-n|)^{-p};
  \]
\item for off-diagonal indices $|m-n|\gg1$, the same algebraic estimate
  holds uniformly away from the diagonal.
\end{enumerate}
Then the quadrature defect satisfies the corresponding partial bound
\[
  |E_{mn}|\le C_p'(1+|m-n|)^{-p}+C e^{-\alpha N}
\]
in the union of these bulk and off-diagonal regimes. In particular, one
obtains a banded-type compatibility statement away from the unresolved edge
region.
\end{proposition}
\begin{proof}
Insert the assumed algebraic tail estimate into
Proposition~\ref{prop:structured_defect}. The $L^2$ term is uniformly
bounded by Lemma~\ref{lem:product_norm}, so the residual contribution
remains exponentially small in $N$. The resulting estimate is therefore of
the stated algebraic-plus-exponential form in every regime where the
assumed partial tail control is available.
\end{proof}

\begin{remark}[Scope of the partial compatibility statement]
Proposition~\ref{prop:partial_compatibility} does not address the full
uniform range $0\le m,n\le N-1$. Its purpose is only to record that bulk and
off-diagonal control would already imply a partial banded structure for the
quadrature defect, while the unresolved edge regime remains the obstruction
to a global compatibility theorem.
\end{remark}

\begin{lemma}[Exponential domination of the bulk error scale]
\label{lem:bulk_scale}
Assume $N=\gamma c$ with fixed $\gamma>0$. Then for every $\alpha>0$ there
exists $C_{\alpha,\gamma}>0$ such that
\[
  e^{-\alpha N}\le C_{\alpha,\gamma} c^{-1/3}
\]
for all sufficiently large $c$.
\end{lemma}
\begin{proof}
Since $N=\gamma c$, one has
$e^{-\alpha N}=e^{-\alpha\gamma c}$. Because exponential decay dominates
any negative power of $c$ as $c\to\infty$, the stated bound follows.
\end{proof}

\begin{remark}[Uniform domination over the index range]\label{rem:index_domination}
For $0\le m,n\le N-1$ one has $|m-n|\le N-1$. Hence any comparison between
an $N$-dependent exponential quantity and an algebraic factor
$(1+|m-n|)^{-p}$ is understood in the worst-case sense over the admissible
index range, namely via $(1+|m-n|)^{-p}\ge (1+N)^{-p}$. In particular,
for every $p\ge1$ and sufficiently large $N$, one has
$e^{-\alpha N}\le C_p(1+N)^{-p}$. This is the precise sense in which an
exponential bound in $N$ dominates the algebraic defect profile required in
Assumption~3.10 of \cite{PaperI}.
\end{remark}

\begin{theorem}[Conditional verification theorem]\label{thm:compatibility_verified}
Assume Conjectures~\ref{conj:pswf_product_tail}
and~\ref{conj:xry_stability}, together with Assumptions~\ref{ass:pswf_sup},
\ref{ass:node_sep}, and~\ref{ass:weights}. Let $(\mathbf{p},\mathbf{w})$
be the XRY PSWF-Gauss rule of order $N=\gamma c$ with fixed $\gamma>0$.
Then for every $p\ge1$ there exist constants $C_p,C>0$ such that
\[
  |E_{mn}|\le C_p(1+|m-n|)^{-p}+Cc^{-1/3}
  \qquad(0\le m,n\le N-1).
\]
Hence Assumption~3.10 of \cite{PaperI} holds conditionally for PSWF-Gauss
sampling under the stated inputs.
\end{theorem}
\begin{proof}
By Proposition~\ref{prop:defect_bound}, $|E_{mn}|\le Ce^{-\alpha N}$.
By Remark~\ref{rem:index_domination}, the exponential bound dominates any
algebraic factor $(1+|m-n|)^{-p}$ uniformly over the index range after
adjusting constants. By Lemma~\ref{lem:bulk_scale}, one also has
$e^{-\alpha N}\le Cc^{-1/3}$ for fixed $\gamma>0$, $N=\gamma c$, and
sufficiently large $c$. Thus both terms in Assumption~3.10 are dominated.
\end{proof}

\begin{remark}[What has been achieved]
The preceding theorem does not close the construction problem of
\cite{PaperI} unconditionally. What it provides is a pure implication:
once Conjectures~\ref{conj:pswf_product_tail} and~\ref{conj:xry_stability}
are established, the compatibility assumption of \cite{PaperI} follows for
XRY PSWF-Gauss quadrature.
\end{remark}

\begin{corollary}[Conditional defect norm bound]\label{cor:defect_norm}
Assume Conjectures~\ref{conj:pswf_product_tail}
and~\ref{conj:xry_stability}, together with Assumption~\ref{ass:pswf_sup}.
Then there exists $C_*>0$ such that
\[
  \norm{\EN}_2\le\norm{\EN}_F\le N C_* e^{-\alpha N}\xrightarrow{N\to\infty}0.
\]
\end{corollary}
\begin{proof}
By Proposition~\ref{prop:defect_bound}, each of the $N^2$ matrix entries of
$\EN$ is bounded by $C_*e^{-\alpha N}$. Hence
$\norm{\EN}_F^2=\sum_{m,n}|E_{mn}|^2\le N^2 C_*^2 e^{-2\alpha N}$, so
$\norm{\EN}_F\le N C_* e^{-\alpha N}$. The operator norm is bounded by the
Frobenius norm, hence $\norm{\EN}_2\le\norm{\EN}_F$.
\end{proof}

%% ==================================================================
\section{Conditional Frame Stability for PSWF-Gauss Sampling}\label{sec:stability}
%% ==================================================================

\begin{theorem}[Conditional frame stability theorem]\label{thm:frame_stability}
Assume Conjectures~\ref{conj:pswf_product_tail}
and~\ref{conj:xry_stability}, together with Assumptions~\ref{ass:pswf_sup},
\ref{ass:node_sep}, and~\ref{ass:weights}. Let $(\mathbf{p},\mathbf{w})$
be the XRY PSWF-Gauss rule of order $N=\gamma c$ with $\gamma<2/\pi$ in the
bulk regime. Once $NC_*e^{-\alpha N}<\lambda_{N-1}(c)$,
$\{(p_j,w_j)\}_{j=1}^N$ is a weighted frame for $V_{N,c}$ with bounds
\[
  A(N,c)\,\norm{f}_{L^2}^2
  \le \sum_{j=1}^N w_j|f(p_j)|^2
  \le B(N,c)\,\norm{f}_{L^2}^2,
  \qquad f\in V_{N,c},
\]
where
\[
  A(N,c)=\lambda_{N-1}(c)-NC_*e^{-\alpha N}>0,
  \qquad
  B(N,c)=\lambda_0(c)+NC_*e^{-\alpha N}.
\]
In the bulk regime, this condition is compatible with the known asymptotic
behaviour of PSWF eigenvalues; see \cite{Slepian1978} and the classical
Landau--Widom theory.
\end{theorem}
\begin{proof}
The Weyl perturbation bound applied to
$\GVM=G^{\mathrm{ref}}_{N,c}+\EN$ gives
$\lambda_{\min}(\GVM)\ge\lambda_{N-1}(c)-\norm{\EN}_2\ge A(N,c)>0$
by Corollary~\ref{cor:defect_norm}.
The upper bound is analogous.
\end{proof}

\begin{remark}[Meaning of the result]
The defect $\EN$ is controlled by two conjectural inputs: the PSWF product
spectral tail conjecture and the XRY stability conjecture. Under these
inputs, the abstract implication framework of \cite{PaperI} becomes
effective for XRY PSWF-Gauss quadrature.
\end{remark}

\begin{remark}[Critical density regime]\label{rem:critical}
The condition $A(N,c)>0$ holds in the Landau--Widom bulk regime
$N<(2/\pi-\varepsilon)c$, where $\lambda_{N-1}(c)\ge1-Ce^{-\eta c}$
dominates the exponential defect. Near the critical density
$N\to(2/\pi)c$, the eigenvalue $\lambda_{N-1}(c)\to0$ and the
stability condition may fail; see Open Problem~\ref{op:critical}.
\end{remark}

%% ==================================================================
\section{Open Problems}\label{sec:open}
%% ==================================================================

\begin{problem}[Derive the PSWF product spectral tail conjecture]
Prove Conjecture~\ref{conj:pswf_product_tail} rigorously, with explicit
constants and explicit dependence on $m$, $n$, $c$, and $T$.
\end{problem}

\begin{problem}[Derive the XRY stability conjecture]
Prove Conjecture~\ref{conj:xry_stability} rigorously from the structural
formulas and numerical analysis of \cite{Xiao2001,OsipovRokhlinXiao},
with explicit constants in terms of $N$, $c$, and the XRY construction
parameters.
\end{problem}

\begin{problem}[Node separation and weight bounds]
Prove Assumptions~\ref{ass:node_sep} and~\ref{ass:weights} quantitatively,
with explicit constants in $N$ and $c$.
\end{problem}

\begin{problem}[Critical density regime]\label{op:critical}
Analyse frame stability for PSWF-Gauss sampling near the critical density
$N\to(2/\pi)c$, where $\lambda_{N-1}(c)\to0$ and the Weyl argument
of Theorem~\ref{thm:frame_stability} degenerates.
\end{problem}

\begin{problem}[Alternative compatible quadrature classes]
Characterise quadrature rules beyond the XRY PSWF-Gauss class satisfying
Assumption~3.10 of \cite{PaperI}.
\end{problem}

%% ==================================================================
\begin{thebibliography}{99}

\bibitem{PaperI}
Anonymous Author,
\emph{Frame Coercivity of Discrete Prolate Sampling Operators
under Quadrature Compatibility: Defect Decomposition and a
Bulk--Turning Decay Mechanism},
preprint, 2026.

\bibitem{OsipovRokhlinXiao}
A.~Osipov, V.~Rokhlin, and H.~Xiao,
\textit{Prolate Spheroidal Wave Functions of Order Zero}, Springer, 2013.

\bibitem{Slepian1978}
D.~Slepian,
\textit{Prolate spheroidal wave functions, Fourier analysis, and
uncertainty~-- V},
Bell Syst.\ Tech.\ J.\ \textbf{57} (1978), 1371--1430.

\bibitem{Xiao2001}
H.~Xiao, V.~Rokhlin, and N.~Yarvin,
\textit{Prolate spheroidal wave functions, quadrature, and interpolation},
Inverse Problems \textbf{17} (2001), 805--838.

\end{thebibliography}

\end{document}
